\begin{textsample}{POS Dim 6 – vem\_ed\_15 – Score 7.00 – vem\_ed\_15/t063.txt}  \label{ex:f6_pos_019}
AOS LEITORES Esta edição está inteiramente dedicada ao IVES Cesu 2022 – International Virtual Exchange Symposium, realizado em 8 de \textbf{dezembro} pela equipe dos Projetos Colaborativos Internacionais (PCIs / Cesu) e transmitido pelo Canal do YouTube da Cesu.
O evento foi organizado em três blocos: a abertura (9h) contou com a participação de representantes da Unidade do Ensino Superior de Graduação do Centro Paula Souza (Cesu – CPS), da Assessoria de Relações Internacionais (ARInter – CPS) e da State University of New York (SUNY / EUA).
No segundo bloco (9h30 às 12h), ocorreram \textbf{apresentações} de parceiros internacionais na realização de PCIs com as Fatecs: Tianjin Normal University / China, Symbiosis Centre for Management Studies (Symbiosis International (Deemed University)/Índia), Durban University of Technology (DUT / África do Sul), SUNY COIL Center (EUA) e Uniminuto (Colômbia).
No período da tarde, foi a vez dos relatos do Centro Paula Souza, iniciando pela \textbf{apresentação} do Eixo de Línguas e Projetos Internacionais.
Depois, houve as \textbf{apresentações} das gestoras pedagógicas regionais Danila Bertolin, Fernanda Coelho e Juliana Tonon.
Isabel Cristina Contro Castaldo, da ARInter, trouxe a experiência do Programa de Aprendizagem Colaborativa Internacional (ProCin), que levou às Etecs, em 2021, a prática de Intercâmbios Virtuais desenvolvida nos PCIs das Fatecs desde 2013.
Por fim, professores das Fatecs Jaboticabal, Garça e Franco da Rocha \textbf{relataram} PCIs em suas \textbf{unidades}.
Boa leitura!
EXPEDIENTE Expediente CPS Diretora-Superintendente: Laura Laganá Vice-Diretora-Superintendente: Emilena Lorenzon Bianco Chefe de Gabinete: Armando Natal Maurício Expediente Cesu Coordenador Técnico: Rafael Ferreira Alves Diretor Acadêmico-Pedagógico: André Luiz Braun Galvão Gestão Educacional: Willian Marcos Muniz Menezes Departamento Administrativo: Elisete Buttignon EDI - Estruturação e Desenvolvimento Instrucional: Thais Lari Braga Cilli Expediente Línguas e Projetos Colaborativos Internacionais - Cesu Coordenação de Línguas e Projetos Internacionais: Mariane Teixeira Coordenação de Projetos Colaborativos Internacionais: Osvaldo Succi Junior Acompanhamento pedagógico PCI: Ana Carolina Freschi, Neusa Haruka Gritti e Regiane Moreira Expediente VEm Corpo editorial: Ana Carolina Freschi, Mariane Teixeira, Neusa Haruka Gritti, Osvaldo Succi Junior e Regiane Moreira Jornalista responsável e Comunicação: Patrícia Patrício - MTb 25.131 Editoração e diagramação: Fábio Gomes da Silva VEm: Virtual Exchange Medium é um informativo com publicação bimestral da Cesu / CEETEPS: Rua dos Andradas, 140 - Santa Efigênia - 01208-000 - São Paulo - SP

% matched lemmas: apresentação, dezembro, relatar, unidade
\end{textsample}
